% Source base: FactstoMemorize.tex, 2026 Science War Edition (2026-09-15)
% This file deliberately omits exhaustive laureate/biography lists from the handbook.

\finalchapter{11}{Mathematical Prizes and Congresses}

\begin{fact}[Science Quiz Prize Triad]
가장 먼저 구분할 세 수학상은
\[
\boxed{\text{Fields Medal}},\qquad
\boxed{\text{Abel Prize}},\qquad
\boxed{\text{Wolf Prize in Mathematics}}.
\]
문제에서는 보통
\[
\text{prize}\longleftrightarrow\text{laureate}
\longleftrightarrow\text{achievement / area}
\]
의 짧은 연결로 나온다.
\end{fact}

\begin{center}
\small
\renewcommand{\arraystretch}{1.22}
\begin{tabularx}{0.97\textwidth}{@{}L{0.19\textwidth}L{0.23\textwidth}Y@{}}
\toprule
Prize & Origin / cycle & 기억할 anchor \\
\midrule
Fields Medal & 1936 first award; ICM에서 보통 4년마다 & 존 찰스 필즈 (John Charles Fields), Canada; age condition; IMU/ICM \\
Abel Prize & 2003 first award; annual & 닐스 헨리크 아벨 (Niels Henrik Abel); Oslo; no age limit \\
Wolf Prize in Mathematics & 1978 first award; Wolf Foundation & 리카르도 울프 (Ricardo Wolf)는 mathematician이 아니었음; Israel; no age limit \\
\bottomrule
\end{tabularx}
\end{center}

\finalsection{Fields Medal and ICM}

\begin{fact}[Fields Medal Basics]
ICM은 \textit{International Congress of Mathematicians}, IMU는 \textit{International Mathematical Union}. Fields Medal은 IMU가 ICM에서 수여한다. 후보자의 $40$번째 생일이 Congress year의 January 1 이전이면 age condition을 만족하지 못한다. 1966년부터 한 ICM에서 최대 네 명에게 수여할 수 있다.
\end{fact}

\begin{center}
\small
\renewcommand{\arraystretch}{1.20}
\begin{tabularx}{0.98\textwidth}{@{}L{0.10\textwidth}L{0.26\textwidth}Y@{}}
\toprule
Year & Name / place & Quiz anchor \\
\midrule
1936 & Oslo; Lars Ahlfors, Jesse Douglas & first Fields Medals \\
1954 & 장피에르 세르 (Jean-Pierre Serre) & youngest Fields Medalist \\
1990 & 에드워드 위튼 (Edward Witten) & first physicist; undergraduate history major \\
1998 & Andrew Wiles special IMU plaque & Fermat's Last Theorem; not a Fields Medal \\
2006 & 그리고리 페렐만 (Grigori Perelman); 테렌스 타오 (Terence Tao) & Perelman declined; Tao: analysis/combinatorics \\
2010 & 세드리크 빌라니 (C'edric Villani) 등 & Villani: Boltzmann equation, optimal transport; Hyderabad ICM \\
2014 & Seoul; 마리암 미르자하니 (Maryam Mirzakhani) & first woman Fields Medalist \\
2018 & Caucher Birkar; Alessio Figalli; Peter Scholze; Akshay Venkatesh & 2018 Fields quartet; Birkar's medal was stolen and replaced \\
2022 & 허준이 (June Huh); 마리나 비아조우스카 (Maryna Viazovska) & Korean connection; Viazovska second woman \\
2026 & Philadelphia; Hong Wang & Hong Wang third woman Fields Medalist \\
\bottomrule
\end{tabularx}
\end{center}

\begin{fact}[Fields Medal Design]
Obverse에는 아르키메데스 (Archimedes)의 옆모습이 있다. Reverse의 sphere-inscribed-in-cylinder는 Archimedes가 자신의 묘비에 새기길 원했던 결과를 떠올리게 한다.
\[
V_{\text{sphere}}:V_{\text{cylinder}}=2:3.
\]
``Fields Medal design $\leftrightarrow$ Archimedes $\leftrightarrow$ sphere/cylinder''를 하나의 association으로 기억한다.
\end{fact}

\finalsection{Abel Prize}

\begin{fact}[Abel Prize Basics]
노르웨이 의회가 2002년에 현대 Abel Prize를 설립했고 2003년에 처음 수여했다. Norwegian Academy of Science and Letters가 매년 수여하며 age restriction이 없다.
\end{fact}

\begin{center}
\small
\renewcommand{\arraystretch}{1.18}
\begin{tabularx}{0.97\textwidth}{@{}L{0.10\textwidth}L{0.31\textwidth}Y@{}}
\toprule
Year & Laureate & Quiz anchor \\
\midrule
2003 & 장피에르 세르 (Jean-Pierre Serre) & first Abel laureate; algebraic topology/geometry/number theory \\
2004 & Michael Atiyah, Isadore Singer & Atiyah--Singer index theorem \\
2005 & 피터 랙스 (Peter D. Lax) & PDE, applied mathematics \\
2016 & 앤드루 와일스 (Andrew Wiles) & Fermat's Last Theorem \\
2017 & Yves Meyer & wavelets, harmonic analysis \\
2019 & Karen Uhlenbeck & first woman Abel laureate; gauge theory/geometric analysis \\
2025 & 가시와라 마사키 (Masaki Kashiwara) & $D$-modules, algebraic analysis, crystal bases \\
2026 & 게르트 팔팅스 (Gerd Faltings) & Mordell conjecture, arithmetic geometry; also 1986 Fields Medal \\
\bottomrule
\end{tabularx}
\end{center}

\begin{fact}[No Nobel Prize in Mathematics]
Nobel의 유언에는 physics, chemistry, physiology or medicine, literature, peace가 명시되었고 mathematics는 포함되지 않았다. ``연애 경쟁 때문에 수학상을 만들지 않았다''는 유명한 이야기는 근거가 없다.
\end{fact}


\finalchapter{12}{Mathematicians and Mathematical Institutions}

\begin{center}
\small
\renewcommand{\arraystretch}{1.19}
\begin{longtable}{@{}L{0.27\textwidth}L{0.65\textwidth}@{}}
\toprule
Mathematician & 가장 강한 quiz anchors \\
\midrule
\endfirsthead
\toprule
Mathematician & 가장 강한 quiz anchors \\
\midrule
\endhead
\bottomrule
\endlastfoot
유클리드 (Euclid) & 《기하학 원론》 (\textit{Elements}); Euclidean geometry; proof tradition. \\
아르키메데스 (Archimedes) & sphere/cylinder $2:3$; Fields Medal design; method of exhaustion. \\
피보나치 (Fibonacci, Leonardo of Pisa) & Fibonacci sequence; \textit{Liber Abaci}; Hindu--Arabic numerals in Europe. \\
르네 데카르트 (Ren\'e Descartes) & Cartesian coordinates; analytic geometry. \\
피에르 드 페르마 (Pierre de Fermat) & Fermat's Last Theorem; Fermat's little theorem; number theory. \\
아이작 뉴턴 (Isaac Newton) & calculus; mechanics; \textit{Principia}. \\
고트프리트 빌헬름 라이프니츠 (Gottfried Wilhelm Leibniz) & calculus notation $\int$, $\dd x$; independent development of calculus. \\
레온하르트 오일러 (Leonhard Euler) & Basel problem; Euler formula; K\"onigsberg bridges; prolific notation. \\
카를 프리드리히 가우스 (Carl Friedrich Gauss) & quadratic reciprocity; 《산술 연구》 (\textit{Disquisitiones Arithmeticae}); number theory. \\
샤를 에르미트 (Charles Hermite) & proved that $e$ is transcendental. \\
자크 아다마르 (Jacques Hadamard) & prime number theorem (independently with de la Vall'ee Poussin); Hadamard matrices/inequality. \\
조충지 (Zu Chongzhi) & ancient Chinese approximation of $\pi$; $355/113$ is a famous accurate rational approximation. \\
게오르크 칸토어 (Georg Cantor) & set theory; diagonal argument; countable vs uncountable infinity; Cantor set. \\
다비트 힐베르트 (David Hilbert) & 23 problems; Hilbert space; foundations; 1900 Paris address. \\
에미 뇌터 (Emmy Noether) & Noether's theorem: symmetry $\leftrightarrow$ conservation law; modern abstract algebra. \\
앨런 튜링 (Alan Turing) & Turing machine; computability; wartime cryptanalysis. \\
노버트 위너 (Norbert Wiener) & Wiener process/Brownian motion; cybernetics. \\
스리니바사 라마누잔 (Srinivasa Ramanujan) & $1729$; number theory, infinite series; collaboration with Hardy. \\
폴 에르되시 (Paul Erd\H{o}s) & enormous collaboration network; Erd\H{o}s number; combinatorics/number theory. \\
에드워드 위튼 (Edward Witten) & 1990 Fields Medal; mathematical physics; first physicist Fields Medalist; history major. \\
그리고리 페렐만 (Grigori Perelman) & Poincar\'e conjecture; Ricci flow; declined 2006 Fields Medal. \\
테렌스 타오 (Terence Tao) & 2006 Fields Medal; harmonic analysis, PDE, additive combinatorics. \\
허준이 (June Huh) & 2022 Fields Medal; Hodge theory in combinatorics/matroids; Korean mathematical community. \\
\end{longtable}
\end{center}

\begin{fact}[A Few Institution Anchors]
\begin{itemize}
  \item \textbf{IMU}: International Mathematical Union; ICM과 Fields Medal의 국제적 운영 주체.
  \item \textbf{CMI}: Clay Mathematics Institute; Millennium Prize Problems.
  \item \textbf{IAS}: Institute for Advanced Study, Princeton; Einstein, von Neumann 등과 연결되는 독립 연구기관.
  \item \textbf{IH\'ES}: Institut des Hautes \`Etudes Scientifiques, France; pure mathematics/theoretical physics.
  \item \textbf{KIAS}: Korea Institute for Advanced Study; 한국의 대표적 기초과학 연구기관.
\end{itemize}
\end{fact}

\begin{fact}[Recognition over Biography]
수학자 문제는 full biography보다 ``이름 $\leftrightarrow$ 정리/분야/일화/상''을 빠르게 맞추는 것이 중요하다. Birth/death year 자체보다 uniquely identifying clue가 강한 경우를 우선한다.
\end{fact}


\finalchapter{13}{Mathematical History, Works, and Anecdotes}

\begin{center}
\small
\renewcommand{\arraystretch}{1.18}
\begin{tabularx}{0.97\textwidth}{@{}L{0.27\textwidth}Y@{}}
\toprule
Work / event & 기억할 연결 \\
\midrule
《기하학 원론》 (\textit{Elements}) & Euclid; axiomatic geometry; 고대 수학의 가장 유명한 저작 중 하나. \\
《산술 연구》 (\textit{Disquisitiones Arithmeticae}) & Gauss; 1801; modern number theory의 landmark. \\
Hilbert's 23 Problems & 1900 Paris ICM; 20세기 수학의 연구 방향에 큰 영향. \\
The Scottish Book & Lw\'ow/Scottish Caf\'e; Banach, Ulam 등 Polish school; problems recorded in a notebook. \\
K\"onigsberg Bridges & Euler; graph theory의 고전적 출발점; 현재 도시명 Kaliningrad, Russia. \\
\bottomrule
\end{tabularx}
\end{center}

\begin{fact}[Hardy--Ramanujan Number]
Hardy가 병문안하며 타고 온 taxi number $1729$가 평범하다고 말하자 Ramanujan은 즉시
\[
1729=1^3+12^3=9^3+10^3
\]
로 두 양의 cubes의 합으로 두 가지 방법으로 표현되는 가장 작은 수라고 지적했다.
\end{fact}

\begin{fact}[George Dantzig's Homework Story]
대학원생 George Dantzig가 늦게 수업에 들어가 칠판의 두 문제를 homework로 착각해 풀었다. 실제로는 당시 유명한 unsolved statistics problems였다. ``불가능하다는 말을 듣지 않아서 풀었다''는 일화로 자주 소개된다.
\end{fact}

\begin{fact}[Witten's Undergraduate Major]
에드워드 위튼 (Edward Witten)은 학부에서 mathematics나 physics가 아니라 \boxed{history}를 전공했다. 이후 mathematical physics의 핵심 인물이 되었고 1990 Fields Medal을 받았다.
\end{fact}

\begin{fact}[Perelman's Refusals]
그리고리 페렐만 (Grigori Perelman)은 Poincar\'e conjecture 해결로 2006 Fields Medal을 받도록 선정되었으나 거부했다. 이후 Clay의 Millennium Prize money도 받지 않았다.
\end{fact}

\begin{fact}[Birkar's Stolen Fields Medal]
Caucher Birkar의 2018 Fields Medal은 Rio de Janeiro ceremony 직후 도난당했다. IMU는 이후 replacement medal을 수여했다.
\end{fact}

\begin{fact}[Archimedes' Tomb Wish]
Archimedes는 sphere와 circumscribed cylinder의 관계를 자신의 가장 자랑스러운 발견 가운데 하나로 여겼고, 묘비에 그 도형을 새기길 원했다고 전해진다. 이 이야기가 Fields Medal reverse와 연결된다.
\end{fact}

\begin{fact}[A Date Chain Worth Remembering]
\[
\begin{aligned}
1936&:\ 	ext{first Fields Medals},\\
1990&:\ 	ext{Witten},\\
2006&:\ 	ext{Perelman refusal / Tao},\\
2014&:\ 	ext{ICM Seoul / Mirzakhani},\\
2022&:\ 	ext{June Huh / Viazovska},\\
2026&:\ 	ext{Philadelphia / Hong Wang}.
\end{aligned}
\]
연도를 단독 암기하기보다 사건과 묶어서 기억한다.
\end{fact}


\finalchapter{14}{Famous Problems and Conjectures}

\begin{center}
\small
\renewcommand{\arraystretch}{1.20}
\begin{longtable}{@{}L{0.25\textwidth}L{0.16\textwidth}L{0.51\textwidth}@{}}
\toprule
Problem & Status & Quiz anchor \\
\midrule
\endfirsthead
\toprule
Problem & Status & Quiz anchor \\
\midrule
\endhead
\bottomrule
\endlastfoot
Riemann Hypothesis & Open & nontrivial zeros of $\zeta(s)$ should have $\RePart(s)=1/2$. \\
Poincar\'e Conjecture & Solved & Perelman; Ricci flow; only Clay-recognized solved Millennium problem at the Sep. 15 cutoff. \\
Fermat's Last Theorem & Solved & Wiles; no nonzero integer solutions to $x^n+y^n=z^n$ for $n>2$. \\
Goldbach Conjecture & Open & every even integer $>2$ is sum of two primes? \\
Twin Prime Conjecture & Open & infinitely many prime pairs differing by $2$? \\
Four Color Theorem & Solved & every planar map can be colored with four colors; famous computer-assisted proof history. \\
Continuum Hypothesis & Independent & G\"odel and Cohen; cannot be proved or disproved from ZFC if ZFC is consistent. \\
Navier--Stokes existence/smoothness & Evaluation pending at cutoff & Sep. 2026 solution announcement + Lean formalization; Clay still evaluating. \\
\end{longtable}
\end{center}

\begin{fact}[Millennium Prize Problems]
Clay의 일곱 Millennium Prize Problems:
\begin{multicols}{2}
\begin{itemize}
  \item P vs NP
  \item Riemann Hypothesis
  \item Yang--Mills and Mass Gap
  \item Navier--Stokes
  \item Birch and Swinnerton-Dyer
  \item Hodge Conjecture
  \item Poincar\'e Conjecture
\end{itemize}
\end{multicols}
September 15, 2026 cutoff에서 Clay가 공식적으로 solved로 분류한 것은 Poincar\'e conjecture 하나였다. Navier--Stokes는 announced solution의 평가가 진행 중이므로 ``공식 해결''로 단정하지 않는다.
\end{fact}

\begin{fact}[Common Distractor Distinctions]
\begin{itemize}
  \item Fermat's Last Theorem은 Millennium Prize Problem이 아니다.
  \item Goldbach와 twin prime conjecture도 Millennium list에 없다.
  \item Four Color Theorem은 이미 solved이다.
  \item Continuum Hypothesis는 단순한 ``open problem''이 아니라 ZFC에서 independent라는 성격이 핵심이다.
  \item Announced proof, preprint, formal verification, peer-reviewed publication, prize body's official recognition은 서로 다른 status이다.
\end{itemize}
\end{fact}
